\section{Постановка задачи}\label{sec:problem}

\textbf{Сцена и объекты.} Сцена $\Omega$ --- случайный элемент неизвестного распределения $\mathcal D$. В сцене есть объекты запретных классов $\mathcal K$: объект $e$ имеет класс $k_e \in \mathcal K$ и след на полу $F_e \subset \mathbb R^2$. Объединение следов класса $k$ обозначим $F_k = \bigcup_{e:\,k_e = k} F_e$.

\textbf{Позы.} Робот проходит траекторию с истинными позами $T_1, \dots, T_q \in SE(2)$. Оценщик позы выдаёт $\hat T_t = D_t T_t$, где $D_t \in SE(2)$ --- накопленная к моменту $t$ ошибка, $D_1 = I$. Робот наземный, поэтому ошибка плоская: сдвиг по полу и поворот вокруг вертикали.

\textbf{Карта.} По наблюдениям и \emph{оценённым} позам строится сетка $\hat M$. Для клетки $j$ и класса $k$ она хранит долю $p_j(k) \in [0, 1]$ наблюдений, поддерживающих $k$. Область класса при пороге $\lambda$:
\begin{equation}
  A_\lambda(k) = \{\, j : j \text{ занята},\ p_j(k) \ge \lambda \,\}.
\end{equation}

\textbf{Система координат планировщика.} Путь строится в момент $q$ в мире, каким его видит робот. Объект, увиденный в момент $t$, стоит на карте с ошибкой $D_t$, а робот сейчас находится с ошибкой $D_q$. Поэтому истинный след в системе планировщика равен $\hat F_k = D_q F_k$. Если ошибка позы постоянна, карта и робот смещены одинаково, и вреда нет; вредна разница ошибок между моментом наблюдения и моментом запроса.

\textbf{Безопасность.} Путь $\gamma\colon [0, 1] \to \mathbb R^2$ безопасен по классу $k$, если $\operatorname{dist}(\gamma(\tau), \hat F_k) \ge d$ при всех $\tau$; здесь $d$ --- безопасное расстояние ($d = 0{,}5$~м).

\textbf{Задача.} Дано $n$ калибровочных сцен, в которых истина известна. Требуется по ним найти для каждого класса запас $r_k \ge 0$, такой что для новой сцены
\begin{equation}
  \mathbb P\bigl(\text{любой путь с } \operatorname{dist}(\gamma, A_\lambda(k)) \ge d + r_k \text{ безопасен по классу } k\bigr) \ge 1 - \alpha,
  \label{eq:goal}
\end{equation}
и запас $r_k$ возможно меньше. Вероятность берётся по калибровочным сценам, новой сцене и реализациям ошибок. Запас сужает свободное пространство, поэтому меньший запас означает больше решаемых задач.

\begin{assumption}[обменяемость]\label{as:exch}
Калибровочные сцены $\Omega_1, \dots, \Omega_n$ и новая сцена $\Omega_{n+1}$ обменяемы: их совместное распределение не меняется при перестановке. Это выполняется, например, если сцены и ошибки в них независимы и одинаково распределены.
\end{assumption}

Других предположений нет: модель ошибок детектора и оценщика позы не нужна.

\textbf{Гипотезы.} Обозначим $\mathrm{Cov}(r)$ левую часть~\eqref{eq:goal} для запаса $r$.
\begin{enumerate}[label=H\arabic*., nosep]
  \item Запас $r^{\mathrm{lab}}_k$, откалиброванный при точной позе ($D_t \equiv I$), даёт $\mathrm{Cov} < 1 - \alpha$, когда ошибка позы сравнима с ошибкой меток.
  \item Запас $r^{\mathrm{pose}}_k$, откалиброванный по эталонным меткам, даёт $\mathrm{Cov} < 1 - \alpha$, пока ошибка позы меньше ошибки меток.
  \item Совместный запас $r^{\mathrm{joint}}_k$ (разд.~\ref{sec:method}) даёт $\mathrm{Cov} \ge 1 - \alpha$ при любом уровне ошибки позы, и $r^{\mathrm{joint}}_k < r^{\mathrm{lab}}_k + r^{\mathrm{pose}}_k$.
  \item Метрика качества карты mIoU слабо связана с безопасностью миссий: ранговая корреляция по сценам близка к нулю.
\end{enumerate}
